% !TeX root = ../教材题库_全部条目_转写工作稿.tex
% Source: XJ-B1-C06-PR-P253-Y0840
% 第6章 / 6.4 用样本估计总体 / 教材 p.253 / PDF 第260页
% OCR未校对：公式与文字均可能有误；未完成数学TeX转写；禁止直接发给学生
% 原题及相邻原解截图（核对用）：../题目原图/XJ-B1-C06-PR-P253-Y0840-01.webp
% 本文件只在转写完成后作为题干使用；OCR 条目尚不是可靠题干。
\begin{bankitem}{p.253 练习 第2题}{第6章 / 6.4 用样本估计总体 / 教材 p.253 / PDF 第260页}{XJ-B1-C06-PR-P253-Y0840}{OCR未校对：公式与文字均可能有误；未完成数学TeX转写；禁止直接发给学生}
\begin{ocrdraft}
[教材p.253 / PDF第260页]
\par
2．某城市计划对居民生活用气（天然气）按年采用三阶式收费：75％的用户
\par
在最低气价一档，18％的用户用气量超出一阶气价的临界值而未超过二阶气价的
\par
临界值，超过一阶临界值的用气量按二阶气价缴费，7％的用户用气量超过二阶
\par
气价的临界值，超过二阶临界值的用气量按三阶气价缴费．为此，
\par
当地燃气公司
\par
调查了100户居民一年的用气量（单位：m），并排序如下：
\par
105
\par
120
\par
140
\par
142
\par
146
\par
巧5
\par
160
\par
165
\par
178
\par
187
\par
192
\par
199
\par
200
\par
206
\par
206
\par
213
\par
220
\par
223
\par
225
\par
230
\par
233
\par
239
\par
241
\par
245
\par
248
\par
249
\par
252
\par
254
\par
256
\par
256
\par
258
\par
260
\par
263
\par
265
\par
266
\par
267
\par
270
\par
271
\par
272
\par
273
\par
275
\par
278
\par
280
\par
283
\par
286
\par
287
\par
290
\par
290
\par
290
\par
299
\par
300
\par
303
\par
304
\par
305
\par
306
\par
308
\par
310
\par
311
\par
313
\par
316
\par
316
\par
318
\par
321
\par
323
\par
325
\par
326
\par
326
\par
327
\par
329
\par
330
\par
332
\par
333
\par
335
\par
336
\par
336
\par
337
\par
338
\par
340
\par
341
\par
385
\par
396
\par
413
\par
420
\par
428
\par
431
\par
443
\par
454
\par
56
\par
60
\par
65
\par
470
\par
475
\par
80
\par
485
\par
90
\par
497
\par
500
\par
510
\par
520
\par
536
\par
（D阶梯气价的临界点如何确定？
\par
（2）若第一档气价为2．90元/m，第二档气价为3，48元/m，第三档气价为
\par
4，35元/m，某户居民今年用气330m，则应缴纳多少燃气费？
\end{ocrdraft}
\end{bankitem}
